Una espira quadrada de 10,0~cm de costat es desplaça a una velocitat constant de $5,00\ \mathrm{cm\,s^{-1}}$ i entra en una regió on hi ha un camp magnètic uniforme, perpendicular al pla del paper i dirigit cap endins. El mòdul del camp magnètic és de 5,00~mT. L'amplària de la regió amb camp magnètic és de 30,0~cm. Es pren com a instant inicial $t = 0\ \mathrm{s}$, el moment en què l'espira comença a entrar a la regió amb camp magnètic.

\figura[0.4]{fig1.png}

\begin{apartat}{1,25}
Calculeu el valor absolut del flux magnètic que travessa l'espira en els instants $t = 1,00\ \mathrm{s}$, $t = 5,00\ \mathrm{s}$ i $t = 7,00\ \mathrm{s}$.
\end{apartat}

\begin{apartat}{1,25}
Representeu a la quadrícula de sota el valor absolut de la força electromotriu induïda en funció del temps durant els primers 9,00~s. Justifiqueu el sentit del corrent induït en els diferents intervals en què es genera la força electromotriu.
\end{apartat}

\figura[0.5]{fig2.png}
